% Part: incompleteness
% Chapter: incompleteness-provability
% Section: provability-conditions

\documentclass[../../../include/open-logic-section]{subfiles}

\begin{document}

\olfileid{inc}{inp}{prc}

\olsection{شروط \usetoken{S}{derivability} لـ$\Th{PA}$}

حساب بيانو، الذي نرمز إليه بـ$\Th{PA}$، امتداد لـ$\Th{Q}$
ببديهيات الاستقراء لجميع الـ!!{formula}s. فنضيف إلى $\Th{Q}$
كل بديهية من الصورة
\[
(!A(\OLMathNumeral{0}) \land \lforall[\OLId{latin-ordinary}{x}][(!A(\OLId{latin-ordinary}{x}) \lif !A(\OLId{latin-ordinary}{x}'))]) \lif \lforall[\OLId{latin-ordinary}{x}][!A(\OLId{latin-ordinary}{x})]
\]
حيث !!{formula} $!A$ كيفما كانت. فهذا \emph{مخطط} لا بديهية
واحدة: تندرج تحته بديهيات لا نهاية لها، بل $\Th{PA}$ {\em لا}
تقبل الأكسمة بمجموعة منتهية. ومع ذلك، يمكن فحص سلسلة الرموز
فحصًا فعالًا لمعرفة أهي حالة من بديهية الاستقراء؛ فتكون مجموعة
بديهيات $\Th{PA}$ قابلة للحساب. و$\Th{PA}$ أقوى كثيرًا من
$\Th{Q}$؛ فمن اليسير، مثلًا، إثبات إبدالية الجمع والضرب فيها
بالاستقراء المعتاد. وأكثر الحجج التناهية في نظرية الأعداد
والتوافقيات يمكن إجراؤه في $\Th{PA}$.

لأكسمة $\Th{PA}$ الحسابية يكون محمول !!{derivability}
$\Prf[\Th{PA}](\OLId{latin-ordinary}{x},\OLId{latin-ordinary}{y})$ قابلًا للحساب، فيقبل التمثيل في $\Th{Q}$،
ومن ثم في $\Th{PA}$. ولتكن $\OPrf[\Th{PA}](\OLId{latin-ordinary}{x},\OLId{latin-ordinary}{y})$، كما سبق،
الصيغة الممثلة له، ولتكن $\OProv[\Th{PA}](\OLId{latin-ordinary}{y})$ هي
$\lexists[\OLId{latin-ordinary}{x}][\OPrf[\Th{PA}](\OLId{latin-ordinary}{x},\OLId{latin-ordinary}{y})]$.
ومعناها الحدسي أن «$\OLId{latin-ordinary}{y}$ !!{derivable} من بديهيات $\Th{PA}$».
إنما تجاوزنا بديهيات $\Th{Q}$ قليلًا لنضمن أن النظرية قادرة
على !!{derive} جملة من الحقائق الأساسية عن محمول
!!{derivability}. وهذه الحقائق هي:
\begin{enumerate}
\item[P1.] إذا كان $\Th{PA}\Proves !A$، فإن
  $\Th{PA}\Proves\OProv[\Th{PA}](\gn{!A})$.
\item[P2.] لجميع الـ!!{formula}s $!A$ و$!B$،
  \[
  \Th{PA} \Proves \OProv[\Th{PA}](\gn{!A \lif !B}) \lif
  (\OProv[\Th{PA}](\gn{!A}) \lif \OProv[\Th{PA}](\gn{!B})).
  \]
\item[P3.] لكل !!{formula}~$!A$،
  \[
  \Th{PA} \Proves \OProv[\Th{PA}](\gn{!A})
  \lif \OProv[\Th{PA}](\gn{\OProv[\Th{PA}](\gn{!A})}).
  \]
\end{enumerate}
والتحقق من هذه الخواص الثلاث يستدعي تعريف الـ!!{formula}
$\OProv[\Th{PA}](\OLId{latin-ordinary}{y})$ تعريفًا دقيقًا، ثم وصف الـ!!{derivation}s
الصورية المناسبة باستعمال بديهيات $\Th{PA}$.
فالشرطان الأول والثاني يسيران، وأكثر العمل في الثالث؛ فتأمل ما
يحتاج إليه ذلك من عمل \dots. وتفصيله طويل قليل الفائدة هنا،
فلنأخذ أن $\Th{PA}$ تستوفي الخواص الثلاث المذكورة.
والاختيار المعقول للصيغة $\OProv[\Th{PA}](\OLId{latin-ordinary}{y})$ يحقق كذلك:
\begin{enumerate}
\item[P4.] إذا كان
  $\Th{PA}\Proves\OProv[\Th{PA}](\gn{!A})$، فإن
  $\Th{PA}\Proves !A$.
\end{enumerate}
غير أننا لا نعتمد هذه الخاصية الأخيرة في البرهان.

\begin{digress}
وقد ترك غودل التفصيل كما نتركه هنا: عرض في خاتمة بحثه سنة \OLClassicalDigits{1931}
مخطط برهان عدم الاكتمال الثاني، ووعد بإيراد التفاصيل في بحث لاحق،
لكنه لم ينجز ذلك الوعد. ذلك أن من فهم الحجة رأى إمكان استيفائها،
فلم يعد عرض تلك التفاصيل لازمًا له.
\end{digress}

\end{document}
